%=====================================================================
\chapter{Dimensionality Reduction and Self-Organizing Maps}\label{ch:pcasom}
%=====================================================================
\locator{4}{Part IV --- Learning Without Full Labels}

\begin{outcomes}
By the end of this chapter you will be able to:
\begin{tight}
  \item \textbf{Explain} why reducing the number of features helps, with reference to
        the curse of dimensionality.
  \item \textbf{Compute} principal components by hand for a small two-dimensional
        dataset, and the variance each explains.
  \item \textbf{Choose} how many components to keep, and \textbf{use} PCA inside a
        pipeline.
  \item \textbf{Describe} the self-organizing map, and \textbf{perform} one training
        step by hand: best-matching unit, neighbourhood and update.
  \item \textbf{Contrast} PCA and SOM as ways of seeing high-dimensional data.
\end{tight}
\end{outcomes}

\begin{prereq}
\chref{math} (vectors, variance, matrices), \chref{knn} (the curse of
dimensionality), \chref{clustering} (centroids and competitive assignment) and
\chref{ann} (units, weights and a learning rate: a self-organizing map is a neural
network, trained without backpropagation).
\end{prereq}

\begin{keyterms}
dimensionality reduction $\bullet$ feature selection vs.\ extraction $\bullet$
covariance matrix $\bullet$ principal component $\bullet$ eigenvector and
eigenvalue $\bullet$ explained variance $\bullet$ scree plot $\bullet$ projection
$\bullet$ reconstruction $\bullet$ self-organizing map (Kohonen map) $\bullet$
competitive learning $\bullet$ best-matching unit $\bullet$ neighbourhood function
$\bullet$ topology preservation $\bullet$ U-matrix
\end{keyterms}

\section{Why Fewer Dimensions}

\chref{knn} showed what happens when there are too many features: distances
concentrate, neighbours stop being meaningful, and every method that relies on
distance degrades. High-dimensional data is also slow to process, impossible to
look at, and often redundant --- the forty metrics a project dashboard records
each week (hours, story points, open bugs, test counts, change requests, \dots) rise
and fall together, because a handful of underlying facts about the project drive
them all. Two responses exist. \term{Feature selection}
keeps a subset of the original features (a tree's splits, lasso's non-zero weights).
\term{Feature extraction} builds a smaller set of new features from all the old
ones. This chapter covers the two extraction methods in the HEC outline: principal
component analysis, which is linear and has an exact solution, and the
self-organizing map, a neural network that produces a two-dimensional map.

\section{Principal Component Analysis}

\begin{definitionbox}[Principal Components]
The first \term{principal component} of a dataset is the direction along which the
data varies most. The second is the direction of greatest remaining variance at right
angles to the first, and so on. Computationally: centre the data, form its
\term{covariance matrix} $S$ ($d \times d$), and find its \term{eigenvectors} ---
directions $\mathbf{v}$ with $S\mathbf{v} = \lambda\mathbf{v}$. Sorted by
\term{eigenvalue} $\lambda$, the eigenvectors are the principal components, and each
$\lambda$ is the variance of the data along its component. The fraction
$\lambda_j / \sum_i \lambda_i$ is the \term{explained variance} of component $j$.
\end{definitionbox}

Keeping the first $m$ components and \term{projecting} each example onto them ---
$z_j = \mathbf{v}_j \cdot (\mathbf{x} - \bar{\mathbf{x}})$ --- replaces $d$ features by
$m$ while losing as little variance as any linear projection can. Equivalently, it is
the projection with the smallest \term{reconstruction} error: mapping the $m$ scores
back into the original space loses the least.

\begin{worked}[PCA by Hand in Two Dimensions]
Five examples: $(1, 2)$, $(2, 1)$, $(3, 3)$, $(4, 5)$, $(5, 4)$.

\smallskip
\textbf{Centre.} The mean is $(3, 3)$; centred data: $(-2, -1)$, $(-1, -2)$, $(0, 0)$,
$(1, 2)$, $(2, 1)$.

\textbf{Covariance} (dividing by $n - 1 = 4$):
\[
\operatorname{var}(x_1) = \tfrac{4 + 1 + 0 + 1 + 4}{4} = 2.5,\qquad
\operatorname{var}(x_2) = 2.5,\qquad
\operatorname{cov}(x_1, x_2) = \tfrac{2 + 2 + 0 + 2 + 2}{4} = 2
\]
\[
S = \begin{pmatrix} 2.5 & 2 \\ 2 & 2.5 \end{pmatrix}
\]
\textbf{Eigenvectors.} For this symmetric pattern they are $(1, 1)/\sqrt{2}$ and
$(1, -1)/\sqrt{2}$. Check: $S(1,1)^{\top} = (4.5, 4.5)^{\top} = 4.5\,(1,1)^{\top}$, and
$S(1,-1)^{\top} = (0.5, -0.5)^{\top} = 0.5\,(1,-1)^{\top}$. So
$\lambda_1 = \mathbf{4.5}$, $\lambda_2 = \mathbf{0.5}$.

\textbf{Explained variance:} PC1 carries $4.5/5 = \mathbf{90\%}$, PC2 the remaining 10\%.

\textbf{Scores on PC1:} $z = (x_1 + x_2)/\sqrt{2}$ for each centred point:
$-2.12, -2.12, 0, 2.12, 2.12$.

\smallskip
One number per example now keeps 90\% of the variation. The discarded direction,
$x_1 - x_2$, only records which of the two coordinates is slightly larger --- if that
is noise, PCA has removed it.
\end{worked}

\dsfig{%
\begin{tikzpicture}[font=\scriptsize, scale=0.95]
\draw[gray!50, -{Stealth[length=1.6mm]}] (0,0) -- (6.3,0) node[right, font=\tiny] {$x_1$};
\draw[gray!50, -{Stealth[length=1.6mm]}] (0,0) -- (0,6.3) node[above, font=\tiny] {$x_2$};
\foreach \t in {1,...,5}{\draw[gray!40] (\t,0.06) -- (\t,-0.06) node[below, font=\tiny] {\t};
                         \draw[gray!40] (0.06,\t) -- (-0.06,\t) node[left, font=\tiny] {\t};}
% PC axes through the mean (3,3)
\draw[greenhl!70!black, line width=1.3pt, -{Stealth[length=2.4mm]}] (3,3) -- (5.12,5.12);
\draw[greenhl!70!black, line width=0.6pt, dashed] (0.6,0.6) -- (3,3);
\node[font=\tiny\bfseries, greenhl!45!black, anchor=south east] at (5.1,5.2) {PC1: $\lambda = 4.5$ (90\%)};
\draw[purplehl, line width=1.3pt, -{Stealth[length=2.4mm]}] (3,3) -- (3.71,2.29);
\draw[purplehl, line width=0.6pt, dashed] (3,3) -- (2.1,3.9);
\node[font=\tiny\bfseries, purplehl, anchor=north west] at (3.75,2.3) {PC2: $\lambda = 0.5$ (10\%)};
% projections
\foreach \x/\y in {1/2,2/1,4/5,5/4}{
  \pgfmathsetmacro{\p}{(\x+\y)/2}
  \draw[accent!60, line width=0.6pt] (\x,\y) -- (\p,\p);
  \fill[accent!60] (\p,\p) circle (1.3pt);}
\foreach \x/\y in {1/2,2/1,3/3,4/5,5/4}{\fill[secondary] (\x,\y) circle (2.2pt);}
\fill[primary] (3,3) circle (1.5pt) node[below right, font=\tiny] {mean};
\node[anchor=north west, align=left, text width=5.6cm, font=\scriptsize, text=gray!55!black]
     at (7.0,5.8)
     {\textbf{Blue:} the five examples.\\[3pt]
      \textbf{Green arrow:} the first principal component, the direction of greatest
      spread; its eigenvalue is the variance along it.\\[3pt]
      \textbf{Red segments:} each example's projection onto PC1. Keeping only the
      projected position (the score) loses the short red distance --- 10\% of the
      variance.\\[3pt]
      \textbf{Purple:} the second component, at right angles.};
\end{tikzpicture}}%
{The worked example. PCA rotates the axes to line up with the directions of greatest
variance; dropping PC2 projects every point onto the green line.}{pca-2d}

\section{How Many Components?}

\dsfig{%
\begin{tikzpicture}
\begin{axis}[width=9.8cm, height=5.0cm, axis lines=left,
  xlabel={number of components kept}, ylabel={cumulative explained variance},
  tick label style={font=\tiny}, label style={font=\scriptsize},
  xmin=0, xmax=41, ymin=0, ymax=1.03]
\addplot[secondary, line width=1.3pt, mark=*, mark size=0.9pt] coordinates
  {(1,0.260) (2,0.446) (3,0.598) (4,0.734) (5,0.750) (6,0.766) (7,0.780) (8,0.793) (9,0.806) (10,0.819) (11,0.831) (12,0.843) (13,0.853) (14,0.864) (15,0.873) (16,0.881) (17,0.890) (18,0.898) (19,0.905) (20,0.912) (21,0.919) (22,0.926) (23,0.932) (24,0.938) (25,0.944) (26,0.949) (27,0.955) (28,0.960) (29,0.965) (30,0.970) (31,0.974) (32,0.978) (33,0.981) (34,0.985) (35,0.988) (36,0.991) (37,0.993) (38,0.996) (39,0.998) (40,1.000)};
\addplot[gray!55, dashed] coordinates {(0,0.9)(41,0.9)};
\addplot[gray!55, dashed] coordinates {(19,0)(19,0.905)};
\node[font=\tiny, accent, anchor=north west, align=left] at (axis cs:19.6,0.87)
     {90\% of the variance\\in 19 of 40 components};
\draw[primary, line width=0.5pt] (axis cs:4,0.734) circle (3pt);
\node[font=\tiny, primary, anchor=north west, align=left] at (axis cs:4.6,0.70)
     {the elbow: 4 components\\already hold 73\%};
\node[font=\tiny, gray!55!black, anchor=south west, align=left] at (axis cs:24,0.08)
     {1{,}000 projects,\\40 standardised weekly metrics};
\end{axis}
\end{tikzpicture}}%
{Cumulative explained variance for forty weekly project metrics. The curve rises steeply
for four components and then flattens: four underlying factors drive almost everything
the dashboard shows, and the rest is noise.}{scree}

Three rules are in use: keep enough components to explain a target fraction of the
variance (90\% or 95\% are common); keep the components before the \emph{elbow} of
the eigenvalue plot (the \term{scree plot}); or --- when PCA is a step before a
supervised model --- treat the number of components as a hyperparameter and choose it
by cross-validation. On the project metrics of this chapter's lab, logistic
regression predicting which projects finish late scores 54\% on one component, 63\%
on two, 89\% on four, and slightly \emph{less} on ten (88\%) or all forty (88\%):
the four components before the elbow carry everything that matters, and the rest
only adds noise. The 90\% rule would have kept nineteen.

\begin{alertbox}[Standardise First, and Remember What PCA Ignores]
PCA finds directions of large \emph{variance}, so a feature measured in large units
dominates the first component purely because of its units. Standardise first, unless
the features share a scale that means something. And PCA is unsupervised: the
direction of greatest variance need not be the direction that separates the classes.
A small, discarded component can be the one that matters.
\end{alertbox}

\section{Self-Organizing Maps}

A \term{self-organizing map} (SOM, or Kohonen map, after Teuvo Kohonen) is a neural
network that learns to lay high-dimensional data out on a two-dimensional grid, so
that similar examples land on nearby grid cells. It is unsupervised, and its output is
a picture: a map of the data that people can read.

\begin{definitionbox}[Self-Organizing Map]
A grid of units (often $10 \times 10$), each holding a weight vector
$\mathbf{w}_u$ of the same dimension as the data. Training repeats, for randomly chosen
examples $\mathbf{x}$:
\begin{tightnum}
  \item \textbf{Compete:} find the \term{best-matching unit} (BMU), the unit whose
        weight vector is closest to $\mathbf{x}$.
  \item \textbf{Cooperate:} every unit $u$ gets a \term{neighbourhood} weight
        $h(u)$, 1 at the BMU and falling with distance \emph{on the grid} ---
        usually $h(u) = \exp\!\big(-\text{grid distance}^2 / 2\sigma^2\big)$.
  \item \textbf{Adapt:} move every unit towards the example:
        $\mathbf{w}_u \leftarrow \mathbf{w}_u + \eta\,h(u)\,(\mathbf{x} - \mathbf{w}_u)$.
\end{tightnum}
Both the learning rate $\eta$ and the neighbourhood width $\sigma$ shrink during
training: first the map organises roughly, then each unit fine-tunes.
\end{definitionbox}

\dsfig{%
\begin{tikzpicture}[font=\scriptsize]
% grid of units
\foreach \r in {0,...,4}{\foreach \c in {0,...,4}{
  \pgfmathsetmacro{\dd}{sqrt((\r-2)^2+(\c-3)^2)}
  \pgfmathsetmacro{\sh}{max(0,75-28*\dd)}
  \node[circle, draw=primary!60, fill=accent!\sh, minimum size=0.55cm] (u\r\c) at (\c*0.9,-\r*0.9) {};}}
\node[circle, draw=accent, line width=1.4pt, fill=accent!80, minimum size=0.55cm] at (3*0.9,-2*0.9) {};
\node[font=\tiny\bfseries, text=accent!80!black, anchor=west] at (4.35,-1.8) {BMU};
\draw[accent!70, -{Stealth[length=1.6mm]}] (4.3,-1.8) -- (2.97,-1.8);
\node[font=\scriptsize\bfseries, text=primary] at (1.8,0.8) {the map: a grid of units};
% input vector
\node[rounded corners=2pt, draw=secondary, fill=secondary!10, minimum width=0.5cm,
      minimum height=2.4cm, align=center, font=\tiny] (x) at (-2.6,-1.8)
      {$x_1$\\[2pt]$x_2$\\[2pt]$\vdots$\\[2pt]$x_d$};
\node[font=\scriptsize\bfseries, text=secondary] at (-2.6,-0.25) {input $\mathbf{x}$};
\foreach \r/\c in {0/0,2/0,4/0,1/1,3/1}{\draw[gray!35] (x.east) -- (u\r\c);}
\node[anchor=north west, align=left, text width=5.6cm, font=\scriptsize, text=gray!55!black]
     at (5.6,0.6)
     {Every unit is connected to every input. The unit whose weights best match
      $\mathbf{x}$ wins (the BMU), and it and its \emph{grid} neighbours --- shaded by
      how strongly they are updated --- all move towards $\mathbf{x}$.\\[4pt]
      Because neighbours are always updated together, neighbouring units end up with
      similar weight vectors: nearby on the map means similar in the data. That is
      \term{topology preservation}.};
\end{tikzpicture}}%
{A self-organizing map. Competition picks one winner; cooperation spreads the update
over its neighbourhood on the grid, which is what makes the map smooth.}{som}

\begin{worked}[One SOM Update]
A one-dimensional map of three units with two-dimensional weights:
$\mathbf{w}_1 = (0.2, 0.8)$, $\mathbf{w}_2 = (0.5, 0.5)$, $\mathbf{w}_3 = (0.9, 0.1)$.
Input $\mathbf{x} = (0.8, 0.3)$; learning rate $\eta = 0.5$; neighbourhood $h = 1$ for
the BMU, $0.5$ for its immediate neighbours, 0 beyond.

\smallskip
\textbf{Compete.} Distances: $\|\mathbf{x} - \mathbf{w}_1\| = \sqrt{0.6^2 + 0.5^2} = 0.781$;
$\|\mathbf{x} - \mathbf{w}_2\| = \sqrt{0.3^2 + 0.2^2} = 0.361$;
$\|\mathbf{x} - \mathbf{w}_3\| = \sqrt{0.1^2 + 0.2^2} = 0.224$. \textbf{BMU: unit 3.}

\textbf{Adapt.}
\begin{tight}
  \item Unit 3 ($h = 1$): $\mathbf{w}_3 + 0.5 \times 1 \times \big((0.8, 0.3) - (0.9, 0.1)\big)
        = (0.9, 0.1) + (-0.05, 0.10) = \mathbf{(0.85, 0.20)}$.
  \item Unit 2 ($h = 0.5$, the BMU's neighbour):
        $(0.5, 0.5) + 0.25 \times (0.3, -0.2) = \mathbf{(0.575, 0.45)}$.
  \item Unit 1 ($h = 0$, two steps away on the map): unchanged.
\end{tight}
Unit 2 moved towards $\mathbf{x}$ although it did not win --- that is the cooperation.
After many such steps, unit 3 represents inputs like $\mathbf{x}$, and unit 2 inputs
somewhere between unit 1's and unit 3's.
\end{worked}

\begin{conceptbox}[Reading a Trained Map]
Each unit becomes a prototype for the examples that map to it, so a SOM is also a
clustering (and with $\sigma \to 0$ it is exactly on-line $k$-means). Colour each unit
by the class of the examples it attracts and the classes appear as regions --- the
lab's six-by-six map of past projects puts maintenance work in one corner and the
two similar kinds, new builds and integrations, next to each other elsewhere. The \term{U-matrix} colours each unit by the
average distance to its neighbours' weights: high values mark the borders between
groups, so cluster boundaries become visible walls on the map.
\end{conceptbox}

\rowcolors{2}{white}{lightbg}
\begin{longtable}{@{}L{3.2cm}L{5.8cm}L{6.0cm}@{}}
\toprule
\rowcolor{primary!15}
\textbf{} & \textbf{PCA} & \textbf{Self-organizing map} \\
\midrule
\endhead
Mapping & Linear projection & Non-linear; follows curved structure \\
Solution & Exact, unique (eigenvectors) & Iterative; depends on initialisation and
schedule \\
Output & $m$ new numeric features & A position on a 2-D grid, plus prototypes \\
Best for & Compression, pre-processing, removing correlation & Visual exploration,
clustering, prototype-based summaries \\
New examples & Project with the same components & Map to the BMU \\
\bottomrule
\end{longtable}

\begin{pitfall}
\begin{tight}
  \item \textbf{Running PCA on unscaled features.} The component is then decided by
        units.
  \item \textbf{Fitting PCA on all the data before splitting.} It is a fitted
        transformer: put it in the pipeline (\chref{workflow}).
  \item \textbf{Assuming high variance means high relevance.} PCA does not look at
        the labels.
  \item \textbf{Interpreting components as causes.} A component is a weighted mix
        of every feature; naming it takes care.
  \item \textbf{Keeping the SOM neighbourhood wide to the end.} The map never
        fine-tunes; shrink $\sigma$ and $\eta$ over training.
\end{tight}
\end{pitfall}

%---------------------------------------------------------------------
\begin{lab}[PCA by Hand, PCA at Work, and a Map of Past Projects]
Part 2 makes 1{,}000 projects whose forty weekly metrics are driven by four hidden
factors (size, quality, churn, staffing) plus noise, and asks which finish late.
Part 3 maps 150 projects of three kinds: maintenance, new builds and integrations,
each described by four standardised features.
\begin{lstlisting}[language=Python]
import numpy as np
from sklearn.decomposition import PCA
from sklearn.pipeline import make_pipeline
from sklearn.preprocessing import StandardScaler
from sklearn.linear_model import LogisticRegression
from sklearn.model_selection import cross_val_score

# --- 1. the worked example: PCA by eigen-decomposition -------------------
X = np.array([[1, 2], [2, 1], [3, 3], [4, 5], [5, 4]], dtype=float)
Xc = X - X.mean(axis=0)                     # centre
S = np.cov(Xc, rowvar=False)                # covariance, divides by n - 1
vals, vecs = np.linalg.eigh(S)              # eigenvalues, ascending
order = np.argsort(vals)[::-1]
vals, vecs = vals[order], vecs[:, order]
print("covariance:\n", S)
print("eigenvalues:", vals, " explained:", (vals / vals.sum()).round(3))
print("PC1 direction:", vecs[:, 0].round(3))
print("scores on PC1:", (Xc @ vecs[:, 0]).round(3))

# --- 2. 40 dashboard metrics: how many components? -----------------------
rng = np.random.default_rng(16)
n, d = 1000, 40                          # projects, weekly metrics
Z = rng.normal(size=(n, 4))              # hidden: size, quality, churn, staff
A = rng.normal(size=(4, d))              # how each metric reflects them
Xm = Z @ A + rng.normal(0, 1.0, (n, d))  # what the dashboard records
ym = (Z[:, 0] - Z[:, 1] + 0.5 * Z[:, 2] + rng.normal(0, 0.5, n) > 0)  # late
pca = PCA().fit(StandardScaler().fit_transform(Xm))
cum = np.cumsum(pca.explained_variance_ratio_)
for target in (0.5, 0.8, 0.9, 0.95):
    print(f"{target:.0%} of variance needs {np.argmax(cum >= target) + 1} "
          f"of {d} components")
for k in (1, 2, 4, 10, 40):
    pipe = make_pipeline(StandardScaler(), PCA(k), LogisticRegression())
    acc = cross_val_score(pipe, Xm, ym, cv=5).mean()
    print(f"logistic regression on {k:2d} components: accuracy {acc:.3f}")

# --- 3. a self-organizing map in twenty lines ----------------------------
rng = np.random.default_rng(16)
kinds = np.repeat([0, 1, 2], 50)         # maintenance, build, integration
centre = np.array([[-1.5, -1.2, -1.5, -0.5],     # small, short, little new
                   [0.8, 0.9, 1.2, 0.2],         # new build
                   [0.8, 0.6, 0.3, 1.0]])        # many integrations
Xi = centre[kinds] + rng.normal(0, 0.45, (150, 4))
Xi = StandardScaler().fit_transform(Xi)
rows, cols = 6, 6
grid = np.array([(r, c) for r in range(rows) for c in range(cols)])
rng = np.random.default_rng(0)
W = rng.normal(0, 0.5, (rows * cols, Xi.shape[1]))
steps = 3000
for t in range(steps):
    x = Xi[rng.integers(len(Xi))]
    bmu = np.argmin(((W - x) ** 2).sum(axis=1))   # best-matching unit
    eta = 0.5 * (1 - t / steps)                   # learning rate decays
    sigma = 3.0 * (1 - t / steps) + 0.5           # neighbourhood shrinks
    d2 = ((grid - grid[bmu]) ** 2).sum(axis=1)    # distance on the map
    h = np.exp(-d2 / (2 * sigma ** 2))            # neighbourhood weight
    W += eta * h[:, None] * (x - W)
# label each unit with the kind of project that most often maps to it
hits = {}
for x, y in zip(Xi, kinds):
    u = np.argmin(((W - x) ** 2).sum(axis=1))
    hits.setdefault(u, []).append(y)
names = "mbi"                       # maintenance, build, integration
for r in range(rows):
    print(" ".join(names[np.bincount(hits[r * cols + c]).argmax()]
                   if r * cols + c in hits else "." for c in range(cols)))
\end{lstlisting}
The last block prints the map: \texttt{m}, \texttt{b} and \texttt{i} mark units that
maintenance, new-build and integration projects map to, and \texttt{.} an unused
unit. Rerun it with \texttt{sigma = 0.5} throughout (no shrinking neighbourhood): is
the map still organised?
\end{lab}

%---------------------------------------------------------------------
\begin{checkpoint}
\begin{tightnum}
  \item A covariance matrix has eigenvalues 6, 3 and 1. What fraction of the variance
        do the first two components explain?
  \item Why must features usually be standardised before PCA?
  \item In a SOM, what is the difference between the best-matching unit and its
        neighbourhood, and why are both updated?
\end{tightnum}
\end{checkpoint}

%---------------------------------------------------------------------
\begin{chaptersummary}
\begin{tight}
  \item Reducing dimensions fights the curse of dimensionality, removes redundancy
        and makes data visible; selection keeps features, extraction builds new
        ones.
  \item PCA's components are the eigenvectors of the covariance matrix; each
        eigenvalue is the variance along its component. Keep enough to explain a
        target fraction, or tune the number by cross-validation.
  \item Standardise before PCA, fit it inside the pipeline, and remember it ignores
        the labels.
  \item A SOM finds the best-matching unit for each input and moves it and its grid
        neighbours towards the input; with a shrinking neighbourhood it produces a
        smooth, topology-preserving 2-D map.
  \item PCA is linear, exact and for compression; the SOM is non-linear, iterative
        and for visual exploration.
\end{tight}
\end{chaptersummary}

\begin{reviewq}
  \item \textbf{(MCQ)} The first principal component is the direction of:
        (a)~greatest class separation (b)~greatest variance (c)~smallest variance
        (d)~the first feature.
  \item \textbf{(MCQ)} In a SOM, the neighbourhood function is measured by distance:
        (a)~in the input space (b)~on the grid of units (c)~between classes
        (d)~between eigenvalues.
  \item \textbf{(MCQ)} Principal components are: (a)~correlated (b)~at right angles
        to each other (c)~always three in number (d)~the original features.
  \item \textbf{(Short)} Explain the difference between feature selection and feature
        extraction, with one method of each.
  \item \textbf{(Short)} Why might a discarded low-variance component matter for
        classification?
  \item \textbf{(Short)} What does a U-matrix show?
  \item \textbf{(Applied)} For the points $(1, 3)$, $(3, 1)$, $(4, 4)$, $(5, 7)$ and
        $(7, 5)$, compute the covariance matrix, verify that $(1, 1)$ is an
        eigenvector, and find the fraction of the variance the first principal
        component explains.
  \item \textbf{(Applied)} Continue the SOM worked example with the next input
        $(0.3, 0.7)$, using the updated weights. Which unit wins, and what are the
        three weight vectors afterwards?
\end{reviewq}
